\section{Poynting vector of the emitted radiation}

In SI units, the instantaneous electromagnetic energy-flux density is the
Poynting vector
\begin{align}
  {\bf S}(t,{\bf x})
  =\frac1{\mu_0}{\bf E}(t,{\bf x})\times{\bf B}(t,{\bf x}).
  \label{eq:standard-Poynting-vector}
\end{align}
Using Eq.~\eqref{eq:fields-from-Faraday}, its Cartesian components can be
written directly in terms of the Faraday tensor as
\begin{align}
  S^i(t,{\bf x})
  =-\frac{c}{\mu_0}\sum_{j=1}^{3}
  F^{j0}(t,{\bf x})F^{ij}(t,{\bf x}).
  \label{eq:Poynting-vector-Faraday}
\end{align}

We use the Fourier convention of
Eq.~\eqref{eq:Faraday-Fourier-convention-observables}, with $e^{i\omega t}$
in the forward transform and $e^{-i\omega t}$ in the inverse transform.
Substitution of the inverse transforms into
Eq.~\eqref{eq:Poynting-vector-Faraday} gives
\begin{align}
  S^i(t,{\bf x})
  =-\frac{c}{\mu_0}
  \int_{-\infty}^{\infty}d\omega
  \int_{-\infty}^{\infty}d\omega'\,
  e^{-i(\omega+\omega')t}
  \sum_{j=1}^{3}
  F^{j0}(\omega,{\bf x})F^{ij}(\omega',{\bf x}).
  \label{eq:instantaneous-Poynting-vector-Fourier-F}
\end{align}

The Poynting vector integrated over time is obtained from Parseval's identity,
\begin{align}
  &\int_{-\infty}^{\infty}dt\,S^i(t,{\bf x})
  =\int_{-\infty}^{\infty}d\omega\,
  \frac{d\mathcal S^i}{d\omega}(\omega,{\bf x}),\\
  &\frac{d\mathcal S^i}{d\omega}(\omega,{\bf x})
  =-\frac{2\pi c}{\mu_0}\operatorname{Re}
  \left[
    \sum_{j=1}^{3}F^{j0}(\omega,{\bf x})
    F^{ij}(\omega,{\bf x})^*
  \right].
  \label{eq:two-sided-spectral-Poynting-vector-F}
\end{align}
Here $d{\boldsymbol{\mathcal S}}/d\omega$ is the two-sided spectral energy
flux density.
\begin{highlightbox}[title={One-sided spectral Poynting vector}]
  For a real time-domain Faraday tensor, the one-sided spectrum
  is
  \begin{align}
    \frac{d\mathcal S^i_+}{d\omega}(\omega,{\bf x})
    =-\frac{4\pi c}{\mu_0}\operatorname{Re}
    \left[
      \sum_{j=1}^{3}F^{j0}(\omega,{\bf x})
      F^{ij}(\omega,{\bf x})^*
    \right],\qquad \omega>0.
    \label{eq:one-sided-spectral-Poynting-vector-F}
  \end{align}
\end{highlightbox}
